\centerline{\bf \Huge Сложный числовой разнобой}

\begin{flushright}
        \scriptsize{--- Никогда не спрашивайте женщину о её возрасте,\\ мужчину о его зарплате и Алину зачем нужна инверсия}
\end{flushright}

\problem{1}
Пусть $a$, $b$, $c$ — произвольные натуральные числа. Докажите, что числа $a^2+b+c$, $b^2+c+a$, $c^2+a+b$ не могут одновременно быть точными квадратами.

\textbf{Решение.}
Без ограничения общности можно считать, что $a\leq b\leq c$. Тогда $c^2<c^2+a+b<c^2+c+c < (c+1)^2$. Значит, число $c^2+a+b$ не может быть точным квадратом.

\problem{2}
Докажите, что из любых десяти последовательных натуральных чисел можно выбрать число, взаимно простое с остальными девятью. Докажите, что произведение десяти последовательных натуральных чисел не может быть точным квадратом.

\textbf{Решение.}
Среди десяти последовательных чисел не более четырех кратны 5 и 7 и не более одного не кратно никакому простому меньшему 10, поэтому минимум пять чисел должны иметь вид $2^a3^bc^2$ где $c$ — некоторое натуральное число. Среди пяти чисел по принципу Дирихле у каких-то двух совпадает четность пары $(a,b)$. Далее рассмотрим четыре случая.

a. Есть два числа с четными показателями $a$ и $b$, то у нас есть два полных квадрата, разность между которыми не превосходит 9. Значит, меньший из них должен быть равен 4, 9 или 16. Во всех случаях нужного произведения десяти последовательных чисел не существует.

b. Есть пара чисел с четным $a$ и нечетным $b$. Тогда есть два утроенных квадрата. Так как разность между ними также не больше 9, то меньший из них должен быть равен 3. Проверкой убеждаемся, что этот случай не подходит.

c. Есть пара чисел с нечетным $a$ и четным $b$. То есть два удвоенных квадрата. Следовательно, меньшее равно 2. Этот вариант также не подходит.

d. Есть пара с нечетными $a$ и $b$, то есть два числа вида $6k^2$. Таких чисел среди десяти подряд идущих не существует.

\problem{3}
Натуральные числа $a$, $b$ и $c$ таковы, что $ab/(a-b)=c$. Известно также, что числа $a$, $b$ и $c$ не имеют общего для них всех натурального делителя, большего 1. Докажите, что $a-b$ есть точный квадрат.

\textbf{Решение.}


\problem{4}
Пусть $M$ — некоторое непустое (возможно, бесконечное) подмножество простых чисел. Известно, что для любого подмножества $K\subset M$ число $(\prod_{p\in K} p )- 1$ раскладывается на простые множители, принадлежащие $M$. Докажите, что тогда $M$ — это все множество простых чисел.

\textbf{Решение.}


\problem{5}
Вася выписал в строку 100 последовательных чисел. Во второй строке под каждым числом он написал его собственный делитель. В третьей строке он под каждым числом из второй строки написал его собственный делитель и т.д., пока не получилось 2022 строк. Могут ли в каждой строке стоять последовательные числа?

\textbf{Решение.}
Докажем, что такое возможно. Для этого начнем строить пример такой таблицы, начиная с последней строки. Поместим в нее числа 2, 3, …, 101. В предпоследнюю строку поместим числа $101!+2$, $101!+3$, …, $101!+101$. Продлим процесс: имея строку $a+1$, $a+2$, …, $a+100$, над ней напишем строку $(a+100)!+a+1$, $(a+100)!+a+2$, …, $(a+100)!+a+100$. Таким образом, мы в конце концов придем к самой первой строке.


\problem{6}
Пусть $p > 5$ — простое число и $S:=\{p - n^2 \colon n \in \mathbb{N}, \, n^2 < p\}$. Докажите, что во множестве $S$ найдутся два числа, одно из которых делит другое.

\textbf{Решение.}
Возьмем такое наибольшее натуральное $n$, что $n^2 < p < (n + 1)^2$.

Рассмотрим случай, когда $p - n^2 = 1$. Тогда $n$ четно. Заметим, что
число $p - 1^2$ делится на $p - (n - 1)^2$:
$$p - (n - 1)^2 = n^2 + 1 - (n - 1)^2 = 2n\mid n^2=p-1^2,$$ При этом
$p - (n - 1)^2 = 2n < n^2 = p - 1$, поскольку в противном случае
$p \leqslant 5$, так что выбранные нами числа различны.

Теперь рассмотрим случай, когда $p - n^2 > 1$. Заметим, что в таком
случае $n \neq p - n^2$, т.к. в противном случае для некоторого простого
$p > 5$ и натурального $n$ верны равенства $n = p - n^2$ и
$p = n^2 + n = n(n + 1)$, что невозможно для простого $p$.

Рассмотрим произвольное $m$ такое, что $m^2 < p$. Тогда $p - m^2$
делится на $p - n^2$ тогда и только тогда, когда разность
$(p - m^2) - (p - n^2)=n^2-m^2$ делится на $p - n^2$.

Положим $m= | n^2 + n - p |$. Ясно, что $m$ — натуральное, причем
$m\neq n$, иначе $p=n^2$ или $p=n^2+2n$, что невозможно. Обозначим
$k = n^2 + n - p$. Тогда $$n - k = n - (n^2 + n - p) = p - n^2.$$
Заметим, что либо $k =  m$, либо $k = -m$, т.е. либо $n - k = n - m$,
либо $n - k = n + m$. Поэтому в произведении $(n - m)(n + m)$ один из
множителей равен $p - n^2$. Значит, $n^2 - m^2$ делится на $p- n^2$, и
$p - m^2$ делится на $p - n^2$, что и требовалось.

\problem{7}
Назовем натуральное число $N$ \textit{сбалансированным}, если или $N=1$, или $N$ раскладывается в произведение четного количества простых чисел с учетом кратности. Пусть также $P(x)=(x+a)(x+b)$ для некоторых натуральных $a$ и $b$.\ Докажите, что найдутся такие различные $a$ и $b$, что все числа $P(1)$, $P(2)$, …, $P(2022)$ сбалансированы.\ Докажите, что если $P(n)$ сбалансировано при любом натуральном $n$, то $a = b$.

\textbf{Решение.}


\newpage

\hspace{3mm}

\problem{8}
На доске написаны числа 1, 2, 3, 4, …, 2022. Петя и Вася играют в игру. Они по очереди стирают написанные числа. Начинает Петя. Тот из мальчиков, после хода которого все оставшиеся числа (быть может, одно) будут иметь натуральный общий делитель, больший единицы, проигрывает. Если же на доске останется только единица, то игра считается завершившейся вничью. Каков будет результат при правильной игре обоих соперников?

\textbf{Решение.}


\problem{9}
Найти все тройки натуральных чисел $(a,b,c)$, такие, что $2^a+2^b+1$ делится на $2^c+1$.

\textbf{Решение.}
Пусть $a_1\equiv a\pmod{c}$ и $b_1\equiv b\pmod{c}$ — остатки чисел $a$ и $b$ при делении на $c$, т.е. $0\leqslant a_1,b_1\leqslant c-1$. Тогда 
$$2^a+2^b+1\equiv 2^{a_1}+2^{b_1}+1\leqslant 2^{c-1}+2^{c-1}+1=2^c+1 \ (\ mod\ 2^c - 1)$$

Если $c = 1$, то $a$ и $b$ — любые.

Если $c\geqslant 2$, то $2^c+1<2(2^c-1)$, поэтому
$$2^{a_1}+2^{b_1}+1=2^c-1\quad\Rightarrow\quad (2^{a_1}-2^{c-1})+(2^{b_1}-2^{c-1})=2.$$

Если обе скобки равны 1, то $c-1=1$, $c=2$, $a_1=b_1=0$ и $a$, $b$ —
любые четные числа.

Если $2^{a_1}-2^{c-1}=2$ и $2^{b_1}-2^{c-1}=0$, то $c-1=2$, $c=3$,
$a_1=1$, $b_1=2$ и $a\equiv 1\pmod{3}$, $b\equiv 2\pmod{3}$.

$c=1$, $a$ и $b$ — любые; $c=2$, $a$ и $b$ — четные; $c=3$,
$a\equiv 1\pmod{3}$, $b\equiv 2\pmod{3}$.

\problem{10}
Для заданного нечетного числа $n>3$ обозначим через $k$ и $\ell$ наименьшие натуральные числа такие, что числа $kn+1$ и $\ell n$ — точные квадраты. Известно, что $k$, $\ell > n/4$. Докажите, что число $n$ --- простое.

\textbf{Решение.}
